\documentclass[11pt]{amsart}
\usepackage[margin=1in]{geometry}
\usepackage{amssymb,mathtools,booktabs}
\usepackage{xurl}
\usepackage[colorlinks=true,linkcolor=blue,citecolor=blue,urlcolor=blue]{hyperref}
\newtheorem{theorem}{Theorem}[section]
\newtheorem{proposition}[theorem]{Proposition}
\newtheorem{lemma}[theorem]{Lemma}
\theoremstyle{remark}
\newtheorem{remark}[theorem]{Remark}
\newcommand{\E}{\mathbb E}

\title[Scope of the sharp contraction bounds]{A remark on the scope of the sharp binary-information\\ and Hellinger contraction bounds}
\author{[AUTHOR NAME]}
\date{October 9, 2026}

\begin{document}
\begin{abstract}
OpenAI's family 119 proves, on the \emph{uniform} binary cube with independent bit-flip noise, a
sharp soft-channel contraction (which implies the Courtade--Kumar conjecture), a refined Boolean
bound depending on the output bias, and the Hellinger conjecture for every output bias. The
statements are made only for the uniform binary cube, and nothing claimed there is contradicted
here. We record that the natural extensions of the refined Boolean bound, of the soft-channel
contraction, and of the Hellinger bound to other finite product spaces (uniform $q$-ary
coordinates with $q\ge3$, or $p$-biased bits with $p\neq\frac12$), \emph{with the standard
resampling noise operator} (the Bonami--Beckner operator; the $q$-ary symmetric channel parametrised
by its correlation), are false, already in dimension one and for small noise. Whether the failure
persists depends on the noise model: with independent bit flips applied to $p$-biased bits, an
exhaustive check in dimensions up to three found no violation. The reason is a second-order computation: on such spaces a single
coordinate can carry a Boolean function of any mean, and its information and Hellinger loss grow
like $\lambda^2/2$ and $\lambda^2/(2\sqrt{1-m^2})$, while the conjectured curves grow like
$\lambda^2r_m^2/2$ and $\lambda^2/2$. The plain Courtade--Kumar curve is not contradicted.
\end{abstract}
\maketitle

\noindent\textbf{Attribution and scope.} The theorems discussed are OpenAI's work: family 119 of
\url{https://github.com/openai/math} (commit \nolinkurl{fd4aeeb2ee4fc729c18d98444fed42fd0529eeeb}),
preprints \cite{oai-ck} and \cite{oai-hel}. Both state their results only for the uniform measure on
$\{-1,1\}^n$; this note concerns only extensions that those preprints do not claim, so it points
out a boundary of applicability, not an error. Developed with AI assistance (Anthropic Claude,
Opus~5.5) under [AUTHOR NAME]'s direction; the independent adversarial review was run by a separate
AI session and is not human peer review.

\section{The statements and their natural extensions}
Logarithms are natural. For $r\in[-1,1]$ let $H(r)$ be the entropy of a sign with mean $r$,
$\psi(r)=\log2-H(r)$, and let $\psi^{-1}$ be its increasing inverse on $[0,1]$. For a Boolean
$f$ of mean $m$ put $r_m=\psi^{-1}(H(m))$, so $r_m<1$ when $m\ne0$. On the uniform cube with
correlation-$\rho$ bit-flip noise, \cite{oai-ck} proves
\begin{equation}\label{eq:soft}
I(T_\rho u)\le\psi\bigl(|\rho|\,\psi^{-1}(I(u))\bigr)\qquad\text{(Theorem 1.1)},
\end{equation}
\begin{equation}\label{eq:refined}
I(f(X);Y)\le\psi(|\rho|\,r_m)\le\psi(|\rho|)\qquad\text{(Corollary 1.2)},
\end{equation}
and \cite{oai-hel} proves
\begin{equation}\label{eq:hel}
\sqrt{1-m^2}-\E\sqrt{1-(T_\rho f)^2}\le1-\sqrt{1-\rho^2}\qquad\text{(Theorem 1.1)}.
\end{equation}
For a Boolean $u=f$ one has $I(f)=H(m)$, so \eqref{eq:soft} specialises to the first inequality
of \eqref{eq:refined}.

\emph{The extension.} Let $(\Omega,\mu)$ be a finite probability space, $X\sim\mu^{\otimes n}$, and
let $R_\lambda$ ($0\le\lambda\le1$) resample each coordinate independently: $Y_i=X_i$ with
probability $\lambda$, otherwise $Y_i$ is a fresh sample of $\mu$. For $\Omega=\{-1,1\}$ with
uniform $\mu$ this is the bit-flip channel with $\rho=\lambda$, and $R_\lambda$ is the standard
noise operator of the $q$-ary and $p$-biased settings (for $q$-ary symmetric noise with symbol
error probability $\varepsilon$, $\lambda=1-\varepsilon q/(q-1)$ is its correlation). Replace
$T_\rho$ by $R_\lambda$ and $|\rho|$ by $\lambda$ in \eqref{eq:soft}--\eqref{eq:hel}.

\section{One-dimensional counterexamples}
Take $n=1$, a set $A\subset\Omega$ with $\mu(A)=w\notin\{0,\tfrac12,1\}$ and $f=2\cdot1_A-1$, so
$m=2w-1\ne0$. Examples: $\Omega=\{1,\dots,q\}$ uniform with $|A|=1$ ($q\ge3$), or a $p$-biased bit
with $A=\{1\}$, $p\ne\frac12$. Since $\E[f(X)\mid Y]=\lambda f(Y)+(1-\lambda)m$,
\[
I(f(X);Y)=H(m)-\E H\bigl(m+\lambda(f-m)\bigr),\qquad
L_{\rm Hel}=\sqrt{1-m^2}-\E\sqrt{1-(m+\lambda(f-m))^2}.
\]

\begin{proposition}\label{prop:main}
As $\lambda\to0$,
\[
I(f(X);Y)=\tfrac{\lambda^2}{2}+O(\lambda^3),\qquad
\psi(\lambda r_m)=\tfrac{\lambda^2r_m^2}{2}+O(\lambda^4),
\]
\[
L_{\rm Hel}=\frac{\lambda^2}{2\sqrt{1-m^2}}+O(\lambda^3),\qquad
1-\sqrt{1-\lambda^2}=\tfrac{\lambda^2}{2}+O(\lambda^4).
\]
Consequently, for $m\ne0$ and all sufficiently small $\lambda>0$, the extensions of
\eqref{eq:refined} (first inequality), of \eqref{eq:soft} with $u=f$, and of \eqref{eq:hel} fail.
\end{proposition}
\begin{proof}
Write $\delta=\lambda(f-m)$; then $\E\delta=0$ and $\E\delta^2=\lambda^2(1-m^2)$. For a smooth
$\Phi$, $\Phi(m)-\E\Phi(m+\delta)=-\tfrac12\Phi''(m)\E\delta^2+O(\lambda^3)$. With
$\Phi=H$, $H''(m)=-1/(1-m^2)$ gives $\lambda^2/2$. With $\Phi(t)=\sqrt{1-t^2}$,
$\Phi''(m)=-(1-m^2)^{-3/2}$ gives $\lambda^2/(2\sqrt{1-m^2})$. Finally $\psi(x)=x^2/2+O(x^4)$, and
$r_m<1$ because $H(m)<\log2$ for $m\ne0$. Compare the coefficients of $\lambda^2$.
\end{proof}

\begin{remark}[mechanism]
For a general concave profile $\Phi$ the same computation gives the coefficient
$\tfrac12(1-m^2)|\Phi''(m)|$, against $\tfrac12|\Phi''(0)|$ for a balanced coordinate. For the
entropy, $(1-m^2)|H''(m)|=1$ for every $m$, so the plain Courtade--Kumar curve $\psi(\lambda)$ is
matched to second order but not exceeded; for $\sqrt{1-t^2}$ the left side is
$(1-m^2)^{-1/2}>1$. On the uniform binary cube a function depending on one coordinate has mean
$0$ or $\pm1$, so no such witness exists there, in agreement with \cite{oai-ck,oai-hel}.
Equivalently: at $n=1$ the whole variance of $f$ lies in the $\lambda$-eigenspace of $R_\lambda$,
whereas on the uniform cube a biased Boolean function cannot put all its variance at level one.
For instance $f=\mathrm{AND}$ on $\{-1,1\}^2$ is measure-isomorphic to the case $q=4$, $|A|=1$,
but under independent bit flips with $\rho=0.1$ it gives $I=0.003345<0.004446$ and Hellinger loss
$0.003864<0.005013$; the violation needs the resampling kernel.
\end{remark}

\begin{remark}[other noise models]
The conclusion depends on the noise model, not only on the space. For the $q$-ary symmetric channel
parametrised by $\rho=1-2\varepsilon$ the violation is larger; parametrised by $\rho=\Pr[Y=X]$ it
disappears (that parametrisation is degenerate at $\rho=0$). For independent bit flips
($Y=X\oplus Z$) applied to $p$-biased bits, the adversarial reviewer checked every Boolean function on
$n\le3$ for $19$ values of $p$ and $25$ values of $\rho$ and found no violation of any of the three
inequalities; at $n=1$ the coefficient of $\rho^2$ is $(1-m^2)/2$, e.g.\ $0.375<r_m^2/2=0.444$ for $w=\tfrac14$.
\end{remark}

\section{Certified witnesses}
Interval arithmetic (mpmath \texttt{iv}, 50 digits, with exact rational inputs) certifies:

\begin{center}
\begin{tabular}{lccccc}
\toprule
setting & $\lambda$ & $I(f;Y)$ & $\psi(\lambda r_m)\le$ & $L_{\rm Hel}$ & $1-\sqrt{1-\lambda^2}$\\
\midrule
$q=4$, $|A|=1$ ($w=\tfrac14$) & $1/10$ & $0.0048170$ & $0.0044458$ & $0.0054645$ & $0.0050126$\\
$q=3$, $|A|=1$ ($w=\tfrac13$) & $1/10$ & $0.0049341$ & $0.0048070$ & $0.0052003$ & $0.0050126$\\
$p$-biased, $p=\tfrac{3}{10}$ & $1/5$ & $0.0193294$ & $0.0188546$ & $0.0207666$ & $0.0202041$\\
\bottomrule
\end{tabular}
\end{center}
In each row both extensions fail, while $I(f;Y)<\psi(\lambda)$. High-precision root finding gives
the violation windows: for $w=\tfrac14$, the refined bound fails for $0<\lambda<0.43648$ and the
Hellinger bound for $0<\lambda<0.32273$; for $w=\tfrac13$, $0<\lambda<0.34953$ and
$0<\lambda<0.33746$.

\section{What is not claimed}
We do not claim that \cite{oai-ck} or \cite{oai-hel} contain an error; both restrict to the uniform
binary cube. The counterexamples concern the resampling noise operator only. We do not claim anything about the plain Courtade--Kumar inequality
$I(f(X);Y)\le\psi(\lambda)$ on general product spaces; our examples respect it, and whether it holds
in that generality is not addressed here. The counterexamples are elementary; we record them only
because the extensions are natural and a reader might expect them to hold.

\begin{thebibliography}{9}
\bibitem{oai-ck} OpenAI, \emph{Sharp binary-information contraction on the discrete cube},
September 24, 2026, \url{https://github.com/openai/math} (family 119).
\bibitem{oai-hel} OpenAI, \emph{Hellinger contraction with arbitrary Boolean output bias},
September 24, 2026, \url{https://github.com/openai/math} (family 119).
\end{thebibliography}
\end{document}
